%%
%% This is file `ic-min-exemplo.tex',
%% generated with the docstrip utility.
%%
%% The original source files were:
%%
%% ufrj-ic.dtx  (with options: `icminexemplo')
%% 
%% This is a sample dissertation or thesis which illustrates the use of the `ufrj'
%% document class with the `ufrj-ic' unit style. It is generated from
%% ufrj-ic.dtx by docstrip.
%% 
%% Copyright (C) 2008-2026 Vicente Helano F. Batista, George O. Ainsworth Jr.,
%% Geraldo Xexéo and Eduardo Mangeli. Released under the GNU General Public
%% License version 3 (see the COPYING.txt file).
%% 
\documentclass[msc]{ufrj}   % msc = dissertacao; dsc = tese
\usepackage{ufrj-ic}
\addbibresource{exemplo.bib}

\begin{document}
  \title{Título do Trabalho}
  \foreigntitle{Title of the Work}
  \author{Nome do Autor}{Sobrenome}
  \advisor{Nome do Orientador}{Sobrenome}{D.Sc.}{UFRJ}
  % A folha de aprovacao lista a banca declarada aqui, na ordem: o orientador,
  % que preside a banca, e o primeiro (3.1.2.1.3e).
  \examiner{Nome do Orientador Sobrenome}{D.Sc.}{PPGI/UFRJ}
  \examiner{Nome do Primeiro Examinador Sobrenome}{D.Sc.}{PPGI/UFRJ}
  \examiner{Nome do Segundo Examinador Sobrenome}{Ph.D.}{UFRJ}
  \department{PPGI}
  \date{03}{2026}
  \approvaldate{30 de março de 2026}
  \keyword{primeira palavra-chave}
  \keyword{segunda palavra-chave}
  \foreignkeyword{first keyword}
  \foreignkeyword{second keyword}

  % Folha adicional da Coleta CAPES (3.1.2.1.2), obrigatoria em dissertacoes e
  % teses. A ficha catalografica vem de fichacatalografica.sibi.ufrj.br; com o
  % arquivo em maos, descomente a linha \catalogcard.
  \concentrationarea{Ciência da Computação}
  \productiontype{bibliographic}
  \linkedproject{no}
  % \catalogcard{ficha.pdf}

  % Ordem dos elementos pre-textuais (Manual 3.1.2), que a classe confere:
  % \maketitle, \frontmatter, o resumo em portugues, o abstract, as listas e
  % \tableofcontents, sempre o ultimo.
  \maketitle
  \frontmatter
  % O resumo e o abstract trazem as instrucoes do Projeto Jacurutu: troque-as
  % pelo seu texto.
  \begin{abstract}
    Utilize, de preferência, um só parágrafo para descrever os objetivos de
    seu trabalho acadêmico, o método com informações sobre coleta, tratamento
    e análise dos dados, as técnicas de abordagem, os resultados e suas
    principais conclusões. Esse resumo deve conter de 150 a 500 palavras. Deve
    ser evitado o uso de frases negativas, símbolos ou contrações que não
    sejam de uso corrente, fórmulas, equações, diagramas etc. que não sejam
    absolutamente necessários; quando for indispensável, defini-los na
    primeira vez que aparecerem. Texto adaptado do Manual do SiBI da UFRJ pelo
    Projeto Jacurutu.
  \end{abstract}
  \begin{foreignabstract}
    Translate your abstract into the language of international dissemination
    chosen by you and your advisor. Translate the keywords as well.
  \end{foreignabstract}
  \tableofcontents

  \mainmatter
  \chapter{Introdução}
  O texto do trabalho começa aqui, e as citações seguem a classe
  \cite{m-4211}.

  \chapter{Conclusão}
  A conclusão do trabalho.

  \backmatter
  \printbibliography
\end{document}
%% 
%%
%% End of file `ic-min-exemplo.tex'.
